% Section 5: Results.
% Every number carries a `% src:` comment naming the committed file it comes
% from (paths relative to the repository root). "md" values are as printed
% there; where a ratio is derived, the comment gives the values it divides.

\section{Results}\label{sec:results}

\subsection{Declared lifetimes are wrong, and a wrong one is worse than none}\label{sec:results-declared}

Table~\ref{tab:declared} compares the lifetime fid's example registry
declares for each information type with the time until half the facts were
wrong.

\begin{table}[t]
\centering
\small
\caption{Declared lifetimes against measured half-lives.}
\label{tab:declared}
\begin{tabular}{llll}
\toprule
Fact & Declared & Measured half-life & Ratio \\
\midrule
Company CEO & 1 year & 5.8 years (empirical) & $5.8\times$ too short \\
Ship position, underway & 4 hours & 3.0 minutes (empirical) & $80\times$ too long \\
Ship position, moored or at anchor & 4 hours & 1.4 days (fitted, no baseline) & $8.4\times$ too short \\
Airport IFR or worse & 6 hours & 5.1 hours (empirical) & $1.2\times$ too long \\
\bottomrule
\end{tabular}
\end{table}
% src: calibration/results/wikidata_ceo.json [0] empirical_half_life_s = 5.788 y; guess.half_life_s = 1 y
% src: calibration/results/ships_ais.json [1] (1 km, underway) empirical_half_life_s = 179.2 s (2.99 min); guess 4 h; 14400 / 179.2 = 80.4
% src: calibration/results/ships_ais.json [2] (1 km, moored) fitted.half_life_s = 1.403 d; 1.403 d / 4 h = 8.42
%      (moored never falls to 0.5 in the 12 h observed, so only a fitted value exists; it is an extrapolation)
% src: calibration/results/weather_metar.json [2] (adverse) empirical_half_life_s = 5.14 h; guess 6 h

The declared value was written for a different purpose in the last row
(weather \emph{forecasts}, not observations), so its closeness is not
evidence that declaring lifetimes works.
% src: docs/FINDINGS.md "Finding 1"

Held out, a declared lifetime is usually worse than ignoring age.
Table~\ref{tab:declared-brier} gives every comparison in the committed
results: the declared lifetime scores worse than the constant in 8 of 10.
These are not independent: six rows are the same ship data at two
tolerances and three groupings. By fact, the declared lifetime loses to the
constant for CEOs, flight category and ship positions, and beats it for
temperature and IFR-or-worse.

\begin{table}[t]
\centering
\small
\caption{Held-out Brier score of the declared lifetime, the constant base
rate, and the fitted model with a baseline (lower is better).}
\label{tab:declared-brier}
\begin{tabular}{lrrr}
\toprule
Fact & Declared & Constant & Fitted, baseline \\
\midrule
Company CEO & 0.3745 & 0.2507 & 0.2008 \\
Airport temperature & 0.2184 & 0.2263 & 0.1809 \\
Airport flight category & 0.4408 & 0.1603 & 0.1554 \\
Airport IFR or worse & 0.2143 & 0.2498 & 0.1990 \\
Ship, 1 km, all & 0.2803 & 0.2256 & 0.2168 \\
Ship, 1 km, underway & 0.4956 & 0.0919 & 0.0299 \\
Ship, 1 km, moored & 0.1707 & 0.0578 & 0.0532 \\
Ship, 10 km, all & 0.1954 & 0.1750 & 0.1529 \\
Ship, 10 km, underway & 0.2437 & 0.2416 & 0.0897 \\
Ship, 10 km, moored & 0.1708 & 0.0384 & 0.0353 \\
\bottomrule
\end{tabular}
\end{table}
% src: calibration/results/wikidata_ceo.md; weather_metar.md (three sections); ships_ais.md (six sections), Brier tables
% Declared worse than constant: CEO, flight category, and all six ship rows (8); better: temperature, IFR (2).
% NOTE: docs/FINDINGS.md says "one scored worse than ignoring age entirely"; the committed results show 8 of 10.

The fitted model with a baseline beats the constant in all ten. For the
CEO fact, the decay-to-zero fit and the fit with a baseline coincide
($b = 0$), with Brier $0.2008$ against $0.2507$ for the constant and
$0.3745$ for the declared year.
% src: calibration/results/wikidata_ceo.md

\subsection{Facts decay toward a baseline, not toward zero}\label{sec:results-baseline}

A week-old report that an airport is VFR is still right most of the time,
because clear weather is the usual state. Decay to zero cannot represent
this; Equation~\eqref{eq:decay} with $b > 0$ can. Table~\ref{tab:baseline}
gives, for every fact, the shape error without and with a baseline, the
fitted baseline, and the held-out Brier of both fits.

\begin{table}[t]
\centering
\small
\caption{Decay to zero against decay to a baseline. Shape error: largest
gap from the empirical curve, all data. Brier: held out.}
\label{tab:baseline}
\begin{tabular}{lrrrrr}
\toprule
 & \multicolumn{2}{c}{Shape error} & & \multicolumn{2}{c}{Brier} \\
\cmidrule(lr){2-3}\cmidrule(lr){5-6}
Fact & to 0 & to $b$ & $b$ & to 0 & to $b$ \\
\midrule
Airport flight category & 0.184 & 0.031 & 0.73 & 0.1664 & 0.1554 \\
Airport IFR or worse & 0.178 & 0.078 & 0.20 & 0.2152 & 0.1990 \\
Airport temperature & 0.415 & 0.183 & 0.25 & 0.2145 & 0.1809 \\
Ship, 1 km, all & 0.314 & 0.067 & 0.62 & 0.2681 & 0.2168 \\
Flight legs & 0.068 & 0.021 & 0.62 & 0.1683 & 0.1682 \\
PyPI latest version, all & 0.102 & 0.056 & 0.15 & 0.1381 & 0.1343 \\
Drug label warnings, all & 0.043 & 0.003 & 0.23 & 0.1515 & 0.1511 \\
Grid demand & 0.621 & 0.292 & 0.33 & 0.3037 & 0.2169 \\
Company CEO & 0.035 & 0.035 & 0.00 & 0.2008 & 0.2008 \\
\bottomrule
\end{tabular}
\end{table}
% src: calibration/results/weather_metar.md, ships_ais.md "AIS 1 km, all vessels", flights_bts.md "Flight legs",
%      pypi_latest.md "all snapshots", dailymed_warnings.md "all snapshots", eia_demand.md, wikidata_ceo.md

Three things follow. First, a baseline costs almost nothing where it is not
needed: where nothing is right by chance, as with a former CEO, the fit
finds $b = 0$ and the two models agree. Across every fact and state in the
committed results, the baseline scores worse held out in only two, both
drug-label states: labels unchanged for over five years ($0.0868$ against
$0.0863$, 36 labels) and labels whose warnings changed in the year before
($0.15272$ against $0.15267$). Second, it matters most where stale facts are often right by
chance: flight category and ship positions. Third, it does not rescue every
fact. Temperature and grid demand still fit poorly (shape error $0.183$
and $0.292$), because both come back daily (Section~\ref{sec:results-echo}).
% src: same rows as Table tab:baseline; "poor fit" threshold 0.10 from calibration/fidcal.py
% src: calibration/results/dailymed_warnings.json brier_test, "unchanged for over 5 years": 0.08633 vs 0.08684;
%      "changed in the year before": 0.15267 vs 0.15272 (equal at 4 decimals in the .md)
%      (checked brier_test in every list-shaped calibration/results/*.json; these are the only two with baseline worse)

When the baseline lies above a decision threshold, the fact never falls
below it, and ``expire after $N$ days'' is the wrong rule: the fact stays
usable at a known, reduced confidence.
% src: docs/FINDINGS.md "Finding 2" (consequence); spec/draft-kumar-fid-00.md "Time to floor" ("If h is infinite, or F <= b', it is never")

\subsection{The half-life depends on the state of the thing observed}\label{sec:results-state}

For the same ships on the same days, position facts recorded while a vessel
was underway were half wrong after 3.0 minutes; those recorded while it was
moored or at anchor, after an extrapolated 1.4 days. After 12 hours, 1.5\%
of underway positions were still within 1\,km, against 79\% of moored ones.
% src: calibration/results/ships_ais.json [1] empirical_half_life_s = 179.2 s; [2] fitted.half_life_s = 1.403 d
% src: calibration/results/ships_ais.md lag 12h rows: underway 0.015 (n = 52,854), moored 0.791 (n = 119,843)

The ratio is about $680$ when the moored fit is compared with the underway
empirical half-life, and about $990$ when both are fitted (the underway fit
is 2.0 minutes). Either way, a single half-life per fact type, the usual
design of expiry policies, is a blend of states that describes neither. For
a moving object the half-life is close to physics, the tolerance divided by
the speed: 1\,km at 12 knots is about 2.7 minutes.
% src: ships_ais.json: 1.4035 d / 179.2 s = 676.8; 1.4035 d / 122.1 s (fitted.half_life_s of [1]) = 993.2
% src: 1 km / (12 kn = 6.173 m/s) = 162 s = 2.7 min (arithmetic)
% NOTE: docs/FINDINGS.md and calibration/README.md say "about 680x" (fitted / empirical) and "held on each of six days";
%       no per-day result is committed, so this draft does not claim the per-day result.

Software shows the same pattern, with a state known when the fact is
stored: how long ago the package's latest release came out
(Table~\ref{tab:pypi}).

\begin{table}[t]
\centering
\small
\caption{``The latest version of package $X$ is $Y$'' by the age of $Y$ when
recorded (PyPI). Fits with a baseline, all data.}
\label{tab:pypi}
\begin{tabular}{lrrrrr}
\toprule
State when recorded & Entities & Half-life & $b$ & Still latest at 1 year & Brier, held out \\
\midrule
released in the 30 days before & 914 & 29 d & 0.07 & 4.8\% & 0.1151 \\
30 days to 1 year old & 886 & 80 d & 0.10 & 15.6\% & 0.1184 \\
no new release for over a year & 496 & 230 d & 0.24 & 50.0\% & 0.1208 \\
all & 915 & 63 d & 0.15 & 17.9\% & 0.1343 \\
\bottomrule
\end{tabular}
\end{table}
% src: calibration/results/pypi_latest.md (four sections); "still latest at 1 year" from pypi_latest.json curve at lag 365d
% Ratio quiet / just released = 230 / 29 = 7.9

Drug label warnings do not show it. Labels whose warnings changed in the
year before and labels unchanged for one to five years have the same
empirical half-life, 2.9 years. Labels unchanged for over five years look
steadier (68\% unchanged after five more years, against 38\% overall), but
there are 36 of them.
% src: calibration/results/dailymed_warnings.md (empirical half-lives 2.9y, 2.9y; entities 36)
% src: calibration/results/dailymed_warnings.json curve at 5.0y: over-5-years 0.678, all 0.380

\subsection{Successive processing steps are not independent}\label{sec:results-steps}

We treat each leg of an aircraft's day as a processing step applied to the
day's published schedule. Across the fleet the on-time rate falls with each
leg flown, from 85.5\% for the first leg to 65.9\% for the seventh, and a
model that multiplies a fixed per-leg loss fits the average: $0.045$ per
leg. It describes no individual aircraft (Table~\ref{tab:legs}).
% src: calibration/results/flights_bts.md "Flight legs" lag table (0 legs 0.855, 6 legs 0.659); "Fitted per-leg noise (fid model): 0.045"

\begin{table}[t]
\centering
\small
\caption{On-time rate of a leg given the previous leg (BTS, 2024).}
\label{tab:legs}
\begin{tabular}{lrr}
\toprule
Previous leg & Next leg on time & Leg pairs \\
\midrule
on time & 85.9\% & 4{,}285{,}796 \\
late & 33.8\% & 1{,}032{,}385 \\
either & 75.8\% & 5{,}318{,}181 \\
\bottomrule
\end{tabular}
\end{table}
% src: calibration/results/flights_bts.md "Flight legs: independence of transform noise"
% "either" n = 4,285,796 + 1,032,385 = 5,318,181

Multiplying per-step loss assumes independence. Where the upstream outcome
is known, the loss should be conditioned on it: after a late leg the step
loses $1 - 0.338/0.758 = 0.55$ of fidelity, not $0.045$.
% src: calibration/results/flights_bts.md ("Conditional noise after a late leg: 1 − 0.338 / 0.758 = 0.55")

The same data shows that a source's reliability is not one number per
class. The published schedule is right (on time within 15 minutes) for
77.9\% of 7{,}079{,}061 flights, but 69.5\% to 84.2\% by airline, 68.1\%
(July) to 85.8\% (October) by month, and 89.6\% for 5\,am departures
against 68.1\% for 8\,pm.
% src: calibration/results/flights_bts.md "Flight schedules: source fidelity" (F9 0.695, YX 0.842; 2024-07 0.681, 2024-10 0.858; 05:00 0.896, 20:00 0.681)

\subsection{Some facts come back}\label{sec:results-echo}

A temperature reading (within $\pm 2\,^\circ$F) is right 19.4\% of the time
at 12 hours old and 43.3\% at 24 hours, because the daily cycle returns it.
No monotonic model can express that. Table~\ref{tab:echo} scores the echo
model of Equation~\eqref{eq:echo} against monotonic models with every hourly
lag to 14 days (29.0 million temperature observations).
% src: calibration/results/weather_periodic.json facts.temperature.curve lag_hours 12 → 0.194, 24 → 0.433
% src: weather_periodic.json facts.temperature.splits.station train_observations + test_observations = 28,999,745

\begin{table}[t]
\centering
\small
\caption{Temperature within $\pm 2\,^\circ$F: held-out Brier in two splits,
and the largest gap from the curve (all data).}
\label{tab:echo}
\begin{tabular}{lrrr}
\toprule
Model & Brier, station split & Brier, time split & Largest gap \\
\midrule
constant & 0.1631 & 0.1746 & 0.776 \\
decay to 0 & 0.2007 & 0.2148 & 0.470 \\
decay to baseline & 0.1599 & 0.1710 & 0.213 \\
echo, $P = 24$\,h & \textbf{0.1572} & \textbf{0.1683} & \textbf{0.070} \\
echo, 24\,h and 168\,h & 0.1572 & 0.1683 & 0.070 \\
\bottomrule
\end{tabular}
\end{table}
% src: calibration/results/weather_periodic.md "METAR temperature"

The echo passes the rule fixed before the run on all three counts. Fitted
on all data, $b = 0.13$, $h = 2.4$\,h, $r = 0.31$, $h_r = 345$\,h and
$w = 3.75$\,h. Scanning the period on the training half alone picks 24\,h,
then 23\,h and 25\,h.
% src: calibration/results/weather_periodic.md (rule lines; "Echo fitted on all data: ... b = 0.128, h = 2.42h, r = 0.311, h_r = 345.0h, w = 3.75h"; period scan)

The Brier gain is small because most observations sit at stale lags, where
every model is near the base rate. The shape is where the echo matters: at
24 hours decay-to-baseline predicts 0.22, the echo 0.39, and the data says
0.43. The model is not right yet. It under-predicts the first echo and
over-predicts the later ones (0.37 against 0.32 at 72 hours), so the first
return is stronger than a single strength fading with $h_r$ allows.
% src: weather_periodic.json facts.temperature.curve: lag 24 → decay_to_baseline 0.220, echo_24h 0.391, still_correct 0.433;
%      lag 72 → echo_24h 0.369, still_correct 0.320

Echo strength varies by station, from $r = 0.17$ at Denver to $0.58$ at Los
Angeles; like the half-life, it belongs to a state.
% src: calibration/results/weather_periodic.md "Temperature echo by station" (DEN 0.173, LAX 0.581)

There is no weekly cycle in weather. Adding a 168-hour echo changes no
held-out Brier at four decimals, and its fitted strength for temperature
is 0.
% src: calibration/results/weather_periodic.md (rows "echo 24h" vs "echo 24h and 168h"; "Weekly term fitted on all data: r7 = 0.0000")

The other two weather facts show the rule's limits. Flight category passes,
but the gain is negligible (station-split Brier $0.19896$ to $0.19887$),
because decay to a $0.71$ baseline already fits it (largest gap $0.045$).
IFR-or-worse fails: the echo improves its Brier but its reliability on the
station split is worse.
% src: weather_periodic.json facts.flight_category.splits.station.brier_test decay_to_baseline 0.198957, echo_24h 0.198869
% src: weather_periodic.md "METAR flight category" (largest gap 0.045; echo b = 0.705) and "METAR adverse conditions" (rule passed: no)
% NOTE: 0.71 is the echo fit's b (0.705); the decay-to-baseline fit in weather_metar.md has b = 0.73 at study lags.

Electricity demand, a fact people drive, echoes too, but not in this
shape. Within $\pm 5\%$, a reading is right 18.9\% of the time at 12 hours,
62.2\% at 24 hours, and 41.7\% at 168 hours against 39.6\% at 144 and 38.3\%
at 192: a weekly echo of two to three points. The echo beats
decay-to-baseline held out (Brier 0.17923 against 0.18430 split by region,
0.16030 against 0.16604 by time), but fails the rule: its largest gap is
0.118, above 0.10, and its reliability is worse in the time split (0.0248
against 0.0219). The first daily return (0.62) is much stronger than the
later ones (about 0.40), which one fading echo cannot express.
% src: calibration/results/eia_demand.json [0].curve lag 12h 0.189, 24h 0.622, 144h 0.396, 168h 0.417, 192h 0.383
% src: calibration/results/eia_demand.md echo test table, rule lines, reliability table

\subsection{Summary of the measurements}

Across seven public sources, the half-life of a stored fact ranged from
minutes (a moving ship) to years (a CEO, a drug label's warnings).
Declared lifetimes were off by up to $80\times$ and usually scored worse
than ignoring age. A baseline, a per-state half-life, conditioned step
noise and, for some facts, a periodic echo each improved held-out
prediction where their effect is present, and cost little or nothing where
it is absent.
% src: sections above
